% ================================================================
% ==== FILE: content/k_levels/k3.tex
% ================================================================

% ==============================
%  Ontology of Continua — Core
%  K-level Module: K3
%  File: content/k_levels/k3.tex
%  Status: FULLY DEFINED
% ==============================

\subsubsection{\texorpdfstring{$K_3$}{K_3} Übersicht}
\label{sec:k3-overview}

$K_3$ ist die Ebene, auf der \emph{chemische Organisation} entsteht
baulich möglich. Es stellt den Übergang von rein geometrisch dar
und topologische Kontinua ($K_0$--$K_2$) zu Systemen, in denen:
\begin{itemize}
    \item Reaktionen, Transformationen und Konvertierungen erscheinen als zulässig
          fließt,
    \item-Potenziale entsprechen Reaktionsaffinitäten und -konzentrationen
          Steigungen,
    \item-Grenzen $\partial\Omega$ umfassen molekular-kompositionelle und
          topologische Einschränkungen,
    \item minimaler autokatalytischer Verschluss wird darstellbar,
    \item entstehen die ersten strukturell zusammenhängenden Netzwerke (RAF-Netzwerke).
\end{itemize}

Der Dimensionsübergang $\Psi_{2\to3}$ tritt auf, wenn er inkompatibel ist
chemische Unterschiede können nicht innerhalb der Achsen $\{A_1,A_2\}$ von ausgedrückt werden
$K_2$. Eine neue Achse $A_3$ wird eingeführt, die eine bestimmte Dimension darstellt
des \emph{Kompositionsraums}, der die Entstehung chemischer Kontinua ermöglicht.

$K_3$ dient als strukturelles Substrat für alle späteren biologischen und
kognitive Kontinua, bleibt aber selbst unbelebt.

% ================================================================
\subsubsection{Zustandsraum \texorpdfstring{$\Omega(K_3)$}{\Omega(K_3)}}

Der Zustandsraum von $K_3$ ist:
\[
\Omega(K_3) = \{ (\mathbf{c}, \mathbf{r}, \Gamma) \mid
\mathbf{c} \in \mathbb{R}_{\ge0}^N,\;
\mathbf{r} \in \mathbb{R}_{\ge0}^M,\;
\Gamma \subseteq R \},
\]
wo:
\begin{itemize}
    \item $\mathbf{c}$: Vektor der Konzentrationen molekularer Spezies,
    \item $\mathbf{r}$: Vektor der Reaktionsflüsse,
    \item $\Gamma$: Satz aktivierter Reaktionen (stöchiometrische Struktur).
\end{itemize}

Eigenschaften:
\begin{enumerate}
    \item $\Omega(K_3)$ ermöglicht nichtlineare Transformationen, die durch induziert werden
          Reaktionsnetzwerke.
    \item Zulässige Zustände müssen Massenbilanzbeschränkungen erfüllen und
          Schwellenwertbeschränkungen ab $M_3$.
    \item RAF-Subnetzwerke sind strukturell in $\Omega(K_3)$ eingebettet
          zusammenhängende Regionen.
\end{enumerate}

Die Stabilität von $K_3$ hängt davon ab, ob $\Omega(K_3)$ geschlossene Werte enthält
autokatalytische Teilmengen mit anhaltendem Fluss.

% ================================================================
\subsubsection{Grenze \texorpdfstring{$\partial\Omega(K_3)$}{\partial\Omega(K_3)}}

$\partial\Omega(K_3)$ beinhaltet:
\begin{itemize}
    \item gibt an, wo jede Konzentration negativ wird (verboten),
    \item divergente Flüsse größer als $\Theta_{\mathrm{flux}}$,
    \item chemische Konfigurationen, die die stöchiometrische Machbarkeit verletzen,
    \item Kollapsregime: pH-Kollaps, osmotischer Kollaps,
          Instabilität der Reaktionskette,
    \item Verlust des autokatalytischen Verschlusses.
\end{itemize}

Kritisch für $K_3$:
\textbf{Verlust der Schließung} impliziert, dass $\Omega(K_3)$ getrennt wird,
eine nachhaltige Organisation verhindern.

% ================================================================
\subsubsection{Achsen \texorpdfstring{$A(K_3)$}{A(K_3)}}

$K_3$ unterstützt drei unabhängige Strukturachsen:
\[
A(K_3) = \{A_1, A_2, A_3\},
\]
wobei $A_3$ die neue Kompositionsachse ist.

$A_3$ kodiert:
\begin{itemize}
    \item stöchiometrische Struktur,
    \item Reaktionsrichtung Freiheitsgrade,
    \item Kompositionsunterschiede, die nicht geometrisch abgebildet werden können
          oder Phasenkohärenzunterschiede von $K_2$.
\end{itemize}

Unabhängigkeit:
\[
A_3 \notin \mathrm{span}\{A_1,A_2\}.
\]

Dies drückt die irreduzible Neuheit der chemischen Organisation aus.

% ================================================================
\subsubsection{Potentiale \texorpdfstring{$P(K_3)$}{P(K_3)}}

Zu den chemischen Potenzialen gehören:
\[
P(K_3) = \{ \mu_i,\; \Delta G_r,\; P_{\mathrm{Affinität}},\;
           P_{\mathrm{mix}},\; P_{\mathrm{Austausch}}\},
\]
mit:
\begin{itemize}
    \item $\mu_i$: Potenziale auf Artenebene,
    \item $\Delta G_r$: freie Gibbs-Reaktionsenergien,
    \item $P_{\mathrm{Affinität}}$: Antriebskräfte der Reaktion,
    \item $P_{\mathrm{mix}}$: entropische Mischungspotentiale,
    \item $P_{\mathrm{exchange}}$: Austausch mit der Umgebung (\emph{if}
          erlaubt durch $M_3$).
\end{itemize}

Chemische Organisation entsteht, wenn diese Potenziale eine stabile,
selbstverstärkende Struktur.

% ================================================================
\subsubsection{Schwellenwerte \texorpdfstring{$\Theta(K_3)$}{\Theta(K_3)}}

Zu den relevanten Schwellenwerten gehören:
\[
\Theta(K_3) =
\{\Theta_{\mathrm{Abschluss}},\;
  \Theta_{\mathrm{Fluss}},\;
  \Theta_{\mathrm{pH}},\;
  \Theta_{\mathrm{osm}},\;
  \Theta_{\mathrm{grad}},\;
  \Theta_{\mathrm{reagieren}}\}.
\]

Interpretation:
\begin{itemize}
    \item $\Theta_{\mathrm{closure}}$: minimale Bedingungen für
          autokatalytischer Verschluss,
    \item $\Theta_{\mathrm{flux}}$: maximale nachhaltige Reaktionsrate,
    \item $\Theta_{\mathrm{pH}}$: zulässiges Säure-Base-Ungleichgewicht,
    \item $\Theta_{\mathrm{osm}}$: Schwelle der osmotischen Stabilität,
    \item $\Theta_{\mathrm{grad}}$: Konzentrations-Gradienten-Grenzen,
    \item $\Theta_{\mathrm{react}}$: Energieschwelle für Reaktionen.
\end{itemize}

Das Überschreiten dieser Schwellenwerte führt zum Zusammenbruch (Abschnitt~the K3 collapse discussion).

% ================================================================
\subsubsection{Flüsse \texorpdfstring{$J(K_3)$}{J(K_3)}}

Zulässige Flüsse:
\[
J(K_3) = \{\text{Reaktionsflüsse},\; \text{diffusiver Transport},\;
         \text{Austauschströme},\; \text{autokatalytische Verstärkung}\}.
\]

Qualitative Typen:
\begin{itemize}
    \item lineare Massenwirkungsflüsse,
    \item nichtlineare autokatalytische Strömungen,
    \item Ausgleichsströme, die eine außer Kontrolle geratene Instabilität verhindern,
    \item-Austausch fließt über Protogrenzen hinweg.
\end{itemize}

Autokatalytische Strömungen sind von zentraler Bedeutung: Sie definieren die strukturelle Verstärkung
notwendig für die Lebensfähigkeit von $K_3$.

% ================================================================
\subsubsection{Zyklen \texorpdfstring{$C(K_3)$}{C(K_3)}}

Zu den Zyklen in $K_3$ gehören:
\[
C(K_3) = \{\text{Reaktionszyklen},\;
           \text{autokatalytische Schleifen},\;
           C_{\mathrm{RAF}}\}.
\]

$C_{\mathrm{RAF}}$ ist der minimale geschlossene RAF-Kreis, der Folgendes erfüllt:
\[
r_i \in \text{RAF} \iff \text{alle Reaktanten von $r_i$ und ein Katalysator
werden innerhalb der Menge erzeugt.}
\]

Eigenschaften:
\begin{itemize}
    \item RAF-Zyklen sorgen für strukturelle Persistenz,
    \item ihr Potenzfunktional $\Pi(C_{\mathrm{RAF}})$ bestimmt
          ob die Organisation nachhaltig ist,
    \item Existenz von $C_{\mathrm{RAF}}$ ist notwendig (aber nicht
          ausreichend) für den Übergang zu $K_4$.
\end{itemize}

Wenn $C_{\mathrm{RAF}}$ stabil wird, unterstützt $K_3$ Vorläuferformen von
zeitliche Ordnung.

% ================================================================
\subsubsection{Zeit \texorpdfstring{$\tau(K_3)$}{\tau(K_3)}}

Zeit in $K_3$ entsteht, wenn autokatalytische Zyklen einen Zyklus ungleich Null aufrechterhalten
Leistung:
\[
\tau(K_3) \text{ existiert genau dann, wenn } \Pi(C_{\mathrm{RAF}}) >
\Theta_{\mathrm{Zeit}}.
\]

Die zeitliche Struktur ist schwächer als in $K_2$:
\begin{itemize}
    Die \item-Zeit wird durch die Iteration des Reaktionszyklus definiert.
    Die \item-Zeit schlägt fehl, wenn die Zyklen erlöschen oder auf den Fluss Null abdriften.
\end{itemize}

Eine stabile zeitliche Ordnung ist eine notwendige Voraussetzung für den Übergang zu
biologisch $K_4$.

% ================================================================
\subsubsection{Kontinuität \texorpdfstring{$k(K_3)$}{k(K_3)}}

Es gilt die allgemeine Formel:
\[
k(K_3) =
\frac{\mu(\Omega(K_3))}{\mu(S(K_3))}
\cdot
\frac{|A(K_3)|}{|A(M_3)|}
\cdot
\frac{\sum \mathrm{Stab}(C)}{\sum \mathrm{MaxStab}(C)}
\cdot
(1 - \frac{\sigma(J)}{\sigma_{\max}})
\cdot
(1 - \frac{T}{\Theta_{\mathrm{stab}}})_+.
\]

Interpretation:
\begin{itemize}
    \item Das Vorhandensein eines RAF-Zyklus steigert $k(K_3)$,
    \item Ein hohes Ungleichgewicht der Flüsse senkt $k(K_3)$,
    \item osmotischer oder pH-Stress reduziert drastisch $k(K_3)$,
    \item $k(K_3)=0$ entspricht dem chemischen Tod.
\end{itemize}

% ================================================================
\subsubsection{Strukturelle Spannung \texorpdfstring{$T(K_3)$}{T(K_3)}}

Spannungsquellen:
\begin{itemize}
    \item inkompatible Reaktionsflüsse,
    \item starke Konzentrationsgradienten,
    \item pH-Ungleichgewicht,
    \item osmotische Gradienten,
    \item fehlende Katalysatoren (unvollständiger Abschluss).
\end{itemize}

Wann:
\[
T(K_3) > \Theta_{\mathrm{stab}},
\]
Die autokatalytische Organisation bricht zusammen.

% ================================================================
\subsubsection{Energie \texorpdfstring{$E(K_3)$}{E(K_3)}}

Energiefunktional:
\[
E(K_3) =
\sum_i \mu_i c_i +
\sum_r \Delta G_r\, r +
E_{\mathrm{grad}} + E_{\mathrm{mix}} + E_{\mathrm{osm}}.
\]

Begriffe bezeichnen:
\begin{itemize}
    \item chemische Potentiale und freie Energien,
    \item Gradienten-Energiekosten,
    \item osmotisches Ungleichgewicht,
    \item Mischentropie.
\end{itemize}

Diese Funktion bestimmt die chemische Stabilität.

% ================================================================
\subsubsection{Operatoren auf \texorpdfstring{$K_3$}{K_3} (\texorpdfstring{$\Psi$}{\Psi}, \texorpdfstring{$\Phi$}{\Phi}, \texorpdfstring{$\Lambda$}{\Lambda}, \texorpdfstring{$U$}{U}, \texorpdfstring{$\Chi$}{\Chi})}

\begin{itemize}
    \item $\Psi_{3\to4}$ — biologischer Übergang: Auftreten von
          Kompartimentierung (Membranverschluss) erzeugt $K_4$.
    \item $\Phi$ – Entwicklung der $M_3$-Einschränkungen (Zusammensetzungsgrenzen,
          Reaktionszulässigkeit).
    \item $\Lambda$ – Operator, der stöchiometrische Konsistenz erzwingt.
    \item $U$ – Vereinheitlichung der Reaktions-Subnetzwerke.
    \item $\Chi$ – ermöglicht die Verzweigung in alternative Chemikalien
          Organisationen (verschiedene RAF-Kerne).
\end{itemize}

% ================================================================
\subsubsection{Prozesse auf \texorpdfstring{$K_3$}{K_3}}

Repräsentative Prozesse:
\begin{itemize}
    \item autokatalytische Verstärkung,
    \item kreuzkatalytische Kettenbildung,
    \item pH-Drift und Erholung,
    \item osmotisches Ungleichgewicht und Stabilisierung,
    \item Bildung oder Auflösung von RAF-Subnetzwerken,
    \item strukturelle Drift in Richtung $K_4$, wenn eine Grenze $\partial\Omega$
          entsteht.
\end{itemize}

% ================================================================
\subsubsection{Vorhersagen für \texorpdfstring{$K_3$}{K_3}}

\begin{enumerate}
    \item Die RAF-Entstehung ist ein Schwellenwertphänomen, das von gesteuert wird
          $\Theta_{\mathrm{Abschluss}}$.
    \item Eine minimale chemische Organisation erfordert autokatalytische Zyklen.
    \item pH-Kollaps führt zu einer schnellen Zerstörung von $\Omega(K_3)$.
    \item Die osmotische Spannungskurve hat eine universelle Sigmoidalform
          ($\Theta_{\mathrm{osm}}$).
    \item Der Übergang zu $K_4$ erfordert die Bildung eines Halbstalls
          Proto-Grenze.
\end{enumerate}

% ================================================================
\subsubsection{Experimente für \texorpdfstring{$K_3$}{K_3}}

Empirische Analoga:
\begin{itemize}
    \item experimentelle RAF-Netzwerke (Steel–Hordijk),
    \item protozellfreie Polymerisationssysteme,
    \item pH-Gradientenoszillatoren,
    \item osmotische Schwellungs-Berst-Zyklen in Vesikeln ohne Membran,
    \item Autokatalytische Reaktoren mit offenem System.
\end{itemize}

Jeder bietet Tests für $\Theta_{\mathrm{closure}}$,
$\Theta_{\mathrm{react}}$ und Kollapskriterien.

% ================================================================
\subsubsection{Zusammenbruch und Tod von \texorpdfstring{$K_3$}{K_3}}
\label{sec:k3-collapse}

Ein Kollaps tritt auf, wenn:
\[
\Omega(K_3) = \varnothing,
\]
aufgrund von:
\begin{itemize}
    \item Zerstörung des RAF-Verschlusses,
    \item pH-Kollaps,
    \item osmotischer Ausbruch,
    \item außer Kontrolle geratene Flussdivergenz,
    \item inkompatible stöchiometrische Einschränkungen.
\end{itemize}

Nach dem Tod ist kein Übergang zum biologischen $K_4$ möglich.

% ================================================================
\subsubsection{Falschbarkeit von \texorpdfstring{$K_3$}{K_3}}

Überprüfbare Vorhersagen:
\begin{enumerate}
    \item Existenz eines minimalen RAF-Schwellenwerts,
    \item Notwendigkeit autokatalytischer Kreisläufe für eine nachhaltige Organisation,
    \item osmotische Kollaps-Hülle reproduzierbar über alle Chemien hinweg,
    \item pH-Grenzwerte, die das strukturelle Überleben definieren,
    \item Unfähigkeit der chemischen Organisation, ohne Schließung zu bestehen.
\end{enumerate}

Verstöße gegen diese Vorhersagen würden das Strukturmodell widerlegen.

% ================================================================
\subsubsection{Verzweigung / ontologische Position von \texorpdfstring{$K_3$}{K_3}}

Die Verzweigung erfolgt durch:
\begin{itemize}
    \item alternative RAF-Kerne,
    \item unterschiedliche stöchiometrische Teilnetze,
    \item Divergente Umweltaustauschregime.
\end{itemize}

Positionell:
\[
K_2 \to K_3 \to K_4,
\]
wobei $K_3$ das einzigartige chemische Zwischenprodukt zwischen physikalischem und ist
biologische Kontinua.

% ================================================================
\subsubsection{Beziehung zu M-Räumen}

$K_3$ existiert innerhalb von $M_3$, was Folgendes angibt:
\begin{itemize}
    \item zulässige Zusammensetzungsbereiche,
    \item zulässige Reaktionsklassen,
    \item Grenzwerte für den Umweltaustausch,
    \item osmotische und pH-Einschränkungen,
    \item strukturelle Bedingungen für die Entstehung von Grenzen $\partial\Omega$.
\end{itemize}

Der Übergang von $K_3 zu K_4$ erfordert, dass $M_3$ einen neuen Zulässigen generiert
Topologie mit einer Innen-/Außenunterscheidung.
